\chapter{Use Cases}

\section*{Motivation}

\begin{remark}
A theoretical modelling framework becomes substantially more convincing
when it is connected to concrete engineering scenarios.
\end{remark}

\begin{remark}
The following use cases illustrate how the proposed AIM framework can be
applied to realistic railway alignment problems.
\end{remark}

\begin{remark}
The examples are not meant as software demonstrations only. They are
intended to clarify the interaction between:
\begin{itemize}
\item data,
\item structured alignment models,
\item dynamic interpretation,
\item and optimization.
\end{itemize}
\end{remark}

\section{Use Case 1: Reconstruction of an Existing Alignment}

\subsection{Scenario}

\begin{remark}
An existing railway alignment is available only through measured points,
legacy file formats, or partially structured engineering data.
\end{remark}

\begin{remark}
Typical input may consist of:
\begin{itemize}
\item a measured polyline,
\item imported TRA/GRA data,
\item LandXML fragments,
\item or mixed data of uncertain quality.
\end{itemize}
\end{remark}

\subsection{Task}

\begin{remark}
The task is to reconstruct a mathematically coherent alignment model from
heterogeneous data.
\end{remark}

\subsection{Pipeline}

\begin{verbatim}
raw data
→ parsing
→ grabbeltisch sorting
→ sparse / pose3 state reconstruction
→ initial guess
→ numerical refinement
\end{verbatim}

\subsection{Interpretation}

\begin{remark}
This use case emphasizes that railway geometry is not simply “read from a
file”. Instead, it is reconstructed as a structured state object.
\end{remark}

\begin{remark}
The AIM framework interprets imported data as evidence for an alignment,
not as unquestionable truth.
\end{remark}

\section{Use Case 2: Smoothing of an Existing Track}

\subsection{Scenario}

\begin{remark}
A measured or existing alignment is geometrically feasible, but exhibits
undesirable local irregularities.
\end{remark}

\begin{remark}
Examples include:
\begin{itemize}
\item noisy curvature estimates,
\item abrupt cant changes,
\item locally degraded transition zones.
\end{itemize}
\end{remark}

\subsection{Task}

\begin{remark}
The task is to derive an improved alignment preserving the overall route
while reducing local dynamic excitation.
\end{remark}

\subsection{Objective Structure}

\begin{remark}
Typical objective terms include:
\begin{itemize}
\item fitting to the original track,
\item smoothness of curvature,
\item smoothness of cant,
\item reduction of jerk.
\end{itemize}
\end{remark}

\subsection{Interpretation}

\begin{remark}
This is one of the clearest examples showing that alignment should be
treated as an optimization problem rather than as a purely geometric
description.
\end{remark}

\section{Use Case 3: Design of a New Transition Zone}

\subsection{Scenario}

\begin{remark}
A new transition between straight track and circular arc is required.
\end{remark}

\begin{remark}
Classically, one would choose a predefined transition family such as a
clothoid. In the present framework, the transition may instead be
derived through optimization.
\end{remark}

\subsection{Task}

\begin{remark}
The task is to construct a transition satisfying:
\begin{itemize}
\item boundary curvature conditions,
\item cant conditions,
\item dynamic smoothness criteria.
\end{itemize}
\end{remark}

\subsection{AIM Perspective}

\begin{remark}
In the AIM framework, the transition is not primarily selected as a
named curve family. Rather, it is derived as a controlled curvature
evolution.
\end{remark}

\begin{remark}
This use case directly illustrates the central thesis that transition
curves are control functions rather than merely geometric connectors.
\end{remark}

\section{Use Case 4: Coupled Optimization of Curvature and Cant}

\subsection{Scenario}

\begin{remark}
A railway section is to be improved with respect to comfort and dynamic
performance.
\end{remark}

\subsection{Task}

\begin{remark}
The task is to optimize both:
\[
\kappa(s)
\quad \text{and} \quad
\phi(s),
\]
instead of treating cant as a secondary afterthought.
\end{remark}

\subsection{Objective Terms}

\begin{remark}
Relevant objective terms include:
\begin{itemize}
\item effective lateral acceleration,
\item jerk,
\item curvature smoothness,
\item cant smoothness.
\end{itemize}
\end{remark}

\subsection{Interpretation}

\begin{remark}
This use case demonstrates the transition from pose2-based geometry to
pose3-based railway state modelling.
\end{remark}

\begin{remark}
It also shows that railway alignment design is inherently a coupled
problem.
\end{remark}

\section{Use Case 5: Schwerpunkt-Trassierung}

\subsection{Scenario}

\begin{remark}
Instead of designing only the track centerline, one seeks to design the
motion of a representative point of the vehicle system.
\end{remark}

\subsection{Task}

\begin{remark}
The task is to optimize the trajectory of the center of mass, or of
another dynamically relevant reference point, by adjusting curvature and
cant.
\end{remark}

\subsection{Model}

\begin{remark}
The design variables are:
\[
\kappa(s), \quad \phi(s),
\]
while the induced trajectory
\[
\gamma_c(s)
\]
becomes the primary object of evaluation.
\end{remark}

\subsection{Interpretation}

\begin{remark}
This use case represents the conceptual extension from classical
alignment design toward the new idea of Schwerpunkt-Trassierung.
\end{remark}

\begin{remark}
It is one of the clearest expressions of the originality of the present
framework.
\end{remark}

\section{Use Case 6: Multi-Track and Network Optimization}

\subsection{Scenario}

\begin{remark}
A railway system rarely consists of a single isolated alignment.
Instead, one must often deal with:
\begin{itemize}
\item parallel tracks,
\item switches,
\item branching structures,
\item coupled corridor geometries.
\end{itemize}
\end{remark}

\subsection{Task}

\begin{remark}
The task is to optimize several alignments simultaneously under
topological and geometric coupling constraints.
\end{remark}

\subsection{Typical Constraints}

\begin{itemize}
\item coincident switch nodes,
\item tangent continuity,
\item parallel track spacing,
\item compatible cant and geometry transitions.
\end{itemize}

\subsection{Interpretation}

\begin{remark}
This use case shows that the AIM framework extends naturally from single
curves to entire railway systems.
\end{remark}

\section{Use Case 7: Interactive Optimization from the Grabbeltisch}

\subsection{Scenario}

\begin{remark}
A user imports multiple heterogeneous data sources and arranges them in
the grabbeltisch interface according to their semantic role.
\end{remark}

\subsection{Task}

\begin{remark}
The system must transform this user-guided sorting into a mathematically
meaningful optimization input.
\end{remark}

\subsection{Pipeline}

\begin{verbatim}
drop files
→ detect / parse
→ sort on grabbeltisch
→ build alignment input
→ construct initial guess
→ optimize
→ preview
→ commit to model
\end{verbatim}

\subsection{Interpretation}

\begin{remark}
This use case is particularly important because it connects the abstract
framework to actual engineering workflow.
\end{remark}

\begin{remark}
It also demonstrates that user interaction and mathematical modelling
are not opposites, but complementary parts of the same system.
\end{remark}

\section{Use Case 8: Standards and BIM Integration}

\subsection{Scenario}

\begin{remark}
An optimized alignment must be exchanged with other systems through
standardized or semi-standardized formats.
\end{remark}

\subsection{Task}

\begin{remark}
The task is to translate the internal alignment state into layered
representations suitable for engineering and BIM workflows.
\end{remark}

\subsection{Interpretation}

\begin{remark}
This use case highlights the value of a structured internal model.
Only such a model can be mapped consistently to multi-layer data models
containing horizontal alignment, profile, and cant.
\end{remark}

\section{Comparative Summary of Use Cases}

\begin{center}
\begin{tabular}{p{4cm}p{4cm}p{5cm}}
\textbf{Use Case} & \textbf{Primary Input} & \textbf{Primary Output} \\
\hline
Reconstruction & raw / imported geometry & structured alignment state \\
Smoothing & measured / degraded track & improved alignment \\
Transition design & boundary conditions & optimized transition law \\
Curvature--cant coupling & route + speed + cant & pose3-consistent state \\
Schwerpunkt & vehicle-based reference & optimized reference trajectory \\
Network optimization & multi-track graph & globally consistent corridor \\
Grabbeltisch workflow & user-sorted imports & optimization-ready input \\
BIM integration & structured state model & exportable layered model \\
\end{tabular}
\end{center}

\section{Conclusion}

\begin{summary}
The presented use cases demonstrate that the AIM framework is not limited
to theoretical alignment descriptions.

It supports:
\begin{itemize}
\item reconstruction,
\item optimization,
\item dynamic interpretation,
\item user-guided engineering workflows,
\item and system-level railway design.
\end{itemize}

Together, these use cases show that alignment can be treated as a
structured, dynamic, and optimization-ready object.
\end{summary}
